%%%%%%%%%%%%%%%%%%%%%%%%%%%%%%%%%%%%%%%%%%%%%%%%%%%%%%%%%%%%
%% Lecture 16
%%%%%%%%%%%%%%%%%%%%%%%%%%%%%%%%%%%%%%%%%%%%%%%%%%%%%%%%%%%%

\lecture
{16}
{Wednesday, 30 September 2026}
{Introduction to Probability}
{Dr. Nishant Chandgotia}


%============================================================
% Strong Law for Bernoulli Random Variables
%============================================================

\section{Strong Law of Large Numbers}

\subsection{Strong Law for Bernoulli Random Variables}

Let
\[
X_1,X_2,\ldots
\overset{\mathrm{i.i.d.}}{\sim}
\operatorname{Ber}\left(\frac12\right),
\]
and define
\[
S_n=X_1+\cdots+X_n.
\]

We want to prove that
\[
\boxed{
\frac{S_n}{n}
\xrightarrow{\mathrm{a.s.}}
\frac12.
}
\]

Fix $\varepsilon>0$ and define the events
\[
A_n
=
\left\{
\left|
\frac{S_n}{n}-\frac12
\right|>\varepsilon
\right\}.
\]

For Bernoulli random variables, Hoeffding's inequality gives the
exponential tail estimate
\[
\mathbb{P}(A_n)
=
\mathbb{P}
\left(
\left|
\frac{S_n}{n}-\frac12
\right|>\varepsilon
\right)
\leq
C'e^{-Cn},
\]
for some constants $C,C'>0$ depending on $\varepsilon$.

Therefore,
\[
\sum_{n=1}^{\infty}\mathbb{P}(A_n)
\leq
C'\sum_{n=1}^{\infty}e^{-Cn}
<\infty.
\]

Hence, by the first Borel--Cantelli lemma,
\[
\boxed{
\mathbb{P}(A_n\text{ i.o.})=0.
}
\]

Equivalently, with probability one, only finitely many of the events
$A_n$ occur.

Thus, for almost every $\omega$, there exists
\[
N=N(\omega,\varepsilon)
\]
such that
\[
\left|
\frac{S_n(\omega)}{n}-\frac12
\right|
\leq\varepsilon
\qquad
\text{for all }n\geq N.
\]

Since $\varepsilon>0$ was arbitrary,
\[
\boxed{
\frac{S_n}{n}
\xrightarrow{\mathrm{a.s.}}
\frac12.
}
\]


%============================================================
% Stronger Tail Estimate
%============================================================

\subsection{A Stronger Tail Estimate}

The above argument can be written in a slightly stronger form.

For every $n$,
\[
\begin{aligned}
&\mathbb{P}
\left(
\left|
\frac{S_k}{k}-\frac12
\right|>\varepsilon
\text{ for some }k>n
\right)
\\
&\qquad=
\mathbb{P}
\left(
\bigcup_{k=n+1}^{\infty}A_k
\right)\\
&\qquad\leq
\sum_{k=n+1}^{\infty}
\mathbb{P}(A_k),
\end{aligned}
\]
where the last inequality follows from the union bound.

Using
\[
\mathbb{P}(A_k)\leq C'e^{-Ck},
\]
we obtain
\[
\begin{aligned}
&\mathbb{P}
\left(
\left|
\frac{S_k}{k}-\frac12
\right|>\varepsilon
\text{ for some }k>n
\right)
\\
&\qquad\leq
C'\sum_{k=n+1}^{\infty}e^{-Ck}\\
&\qquad=
C'
\frac{e^{-C(n+1)}}{1-e^{-C}}.
\end{aligned}
\]

Define
\[
C''
=
\frac{C'e^{-C}}{1-e^{-C}}.
\]

Then
\[
\boxed{
\mathbb{P}
\left(
\left|
\frac{S_k}{k}-\frac12
\right|>\varepsilon
\text{ for some }k>n
\right)
\leq
C''e^{-Cn}.
}
\]

Since
\[
C''e^{-Cn}\longrightarrow0,
\]
we obtain
\[
\boxed{
\mathbb{P}
\left(
\left|
\frac{S_k}{k}-\frac12
\right|>\varepsilon
\text{ for some }k>n
\right)
\longrightarrow0.
}
\]

Now observe that
\[
\begin{aligned}
&
\left\{
\left|
\frac{S_k}{k}-\frac12
\right|>\varepsilon
\text{ infinitely often}
\right\}
\\
&\qquad=
\bigcap_{n=1}^{\infty}
\left\{
\left|
\frac{S_k}{k}-\frac12
\right|>\varepsilon
\text{ for some }k>n
\right\}.
\end{aligned}
\]

The events on the right-hand side form a decreasing sequence.
Therefore, by continuity of probability from above,
\[
\begin{aligned}
&
\mathbb{P}
\left(
\left|
\frac{S_k}{k}-\frac12
\right|>\varepsilon
\text{ i.o.}
\right)
\\
&\qquad=
\lim_{n\to\infty}
\mathbb{P}
\left(
\left|
\frac{S_k}{k}-\frac12
\right|>\varepsilon
\text{ for some }k>n
\right)\\
&\qquad=0.
\end{aligned}
\]

Thus,
\[
\boxed{
\mathbb{P}
\left(
\left|
\frac{S_n}{n}-\frac12
\right|>\varepsilon
\text{ i.o.}
\right)=0.
}
\]

Since this holds for every $\varepsilon>0$, we conclude that
\[
\boxed{
\frac{S_n}{n}
\xrightarrow{\mathrm{a.s.}}
\frac12.
}
\]


%============================================================
% Borel--Cantelli Lemmas
%============================================================

\section{Borel--Cantelli Lemmas}

Let $(A_n)_{n\geq1}$ be a sequence of events.

We say that $A_n$ occurs \emph{infinitely often}, abbreviated as
``$A_n$ i.o.'', if infinitely many of the events $A_n$ occur.

The event that $A_n$ occurs infinitely often is denoted by
\[
\{A_n\text{ i.o.}\}.
\]

We have
\[
\boxed{
\{A_n\text{ i.o.}\}
=
\bigcap_{m=1}^{\infty}
\bigcup_{n=m}^{\infty}A_n.
}
\]

Indeed,
\[
\omega\in\{A_n\text{ i.o.}\}
\]
if and only if, for every $m\geq1$, there exists some $n\geq m$
such that
\[
\omega\in A_n.
\]

Equivalently,
\[
\{A_n\text{ i.o.}\}
=
\limsup_{n\to\infty}A_n.
\]

Similarly,
\[
\boxed{
\{A_n\text{ eventually}\}
=
\bigcup_{m=1}^{\infty}
\bigcap_{n=m}^{\infty}A_n.
}
\]

That is, $A_n$ occurs eventually if there exists some $m$ such that
$A_n$ occurs for every $n\geq m$.


%============================================================
% First Borel--Cantelli
%============================================================

\subsection{First Borel--Cantelli Lemma}

\begin{theorem}{First Borel--Cantelli Lemma}
{first-borel-cantelli}

Let $(A_n)_{n\geq1}$ be a sequence of events. If
\[
\sum_{n=1}^{\infty}\mathbb{P}(A_n)<\infty,
\]
then
\[
\boxed{
\mathbb{P}(A_n\text{ i.o.})=0.
}
\]

\end{theorem}


\begin{proof}

Recall that
\[
\{A_n\text{ i.o.}\}
=
\bigcap_{m=1}^{\infty}
\bigcup_{n=m}^{\infty}A_n.
\]

For each $m\geq1$, define
\[
B_m
=
\bigcup_{n=m}^{\infty}A_n.
\]

Then
\[
B_1\supseteq B_2\supseteq B_3\supseteq\cdots
\]
and
\[
\{A_n\text{ i.o.}\}
=
\bigcap_{m=1}^{\infty}B_m.
\]

By the union bound,
\[
\mathbb{P}(B_m)
=
\mathbb{P}
\left(
\bigcup_{n=m}^{\infty}A_n
\right)
\leq
\sum_{n=m}^{\infty}\mathbb{P}(A_n).
\]

Since
\[
\sum_{n=1}^{\infty}\mathbb{P}(A_n)<\infty,
\]
the tail of this convergent series tends to zero. Hence
\[
\sum_{n=m}^{\infty}\mathbb{P}(A_n)
\longrightarrow0.
\]

Therefore,
\[
\mathbb{P}(B_m)\longrightarrow0.
\]

Since $(B_m)$ is a decreasing sequence of events, continuity of
probability from above gives
\[
\begin{aligned}
\mathbb{P}(A_n\text{ i.o.})
&=
\mathbb{P}
\left(
\bigcap_{m=1}^{\infty}B_m
\right)\\
&=
\lim_{m\to\infty}\mathbb{P}(B_m)\\
&=0.
\end{aligned}
\]

Thus
\[
\boxed{
\mathbb{P}(A_n\text{ i.o.})=0.
}
\]

\end{proof}


%============================================================
% Second Borel--Cantelli
%============================================================

\subsection{Second Borel--Cantelli Lemma}

\begin{theorem}{Second Borel--Cantelli Lemma}
{second-borel-cantelli}

Let $(A_n)_{n\geq1}$ be a sequence of independent events. If
\[
\sum_{n=1}^{\infty}\mathbb{P}(A_n)=\infty,
\]
then
\[
\boxed{
\mathbb{P}(A_n\text{ i.o.})=1.
}
\]

\end{theorem}


\begin{proof}

We prove that the complementary event has probability zero.

The event that only finitely many of the $A_n$ occur is
\[
\{A_n\text{ only finitely often}\}
=
\bigcup_{m=1}^{\infty}
\bigcap_{n=m}^{\infty}A_n^c.
\]

For each $m$, define
\[
C_m
=
\bigcap_{n=m}^{\infty}A_n^c.
\]

Then
\[
C_1\subseteq C_2\subseteq C_3\subseteq\cdots.
\]

Hence, by continuity of probability from below,
\[
\mathbb{P}
\left(
\bigcup_{m=1}^{\infty}C_m
\right)
=
\lim_{m\to\infty}\mathbb{P}(C_m).
\]

Now, for any integer $N\geq m$,
\[
C_m
\subseteq
\bigcap_{n=m}^{N}A_n^c.
\]

Therefore,
\[
\mathbb{P}(C_m)
\leq
\mathbb{P}
\left(
\bigcap_{n=m}^{N}A_n^c
\right).
\]

Since the events $A_n$ are independent, their complements are also
independent. Thus
\[
\begin{aligned}
\mathbb{P}
\left(
\bigcap_{n=m}^{N}A_n^c
\right)
&=
\prod_{n=m}^{N}
\mathbb{P}(A_n^c)\\
&=
\prod_{n=m}^{N}
\left(1-\mathbb{P}(A_n)\right).
\end{aligned}
\]

Using the elementary inequality
\[
1-x\leq e^{-x},
\qquad
0\leq x\leq1,
\]
we obtain
\[
\begin{aligned}
\mathbb{P}(C_m)
&\leq
\prod_{n=m}^{N}
e^{-\mathbb{P}(A_n)}\\
&=
\exp
\left(
-\sum_{n=m}^{N}\mathbb{P}(A_n)
\right).
\end{aligned}
\]

Since
\[
\sum_{n=1}^{\infty}\mathbb{P}(A_n)=\infty,
\]
we have
\[
\sum_{n=m}^{N}\mathbb{P}(A_n)
\longrightarrow\infty
\qquad
\text{as }N\to\infty.
\]

Consequently,
\[
\exp
\left(
-\sum_{n=m}^{N}\mathbb{P}(A_n)
\right)
\longrightarrow0.
\]

Hence
\[
\mathbb{P}(C_m)=0.
\]

Therefore,
\[
\begin{aligned}
\mathbb{P}
(A_n\text{ only finitely often})
&=
\mathbb{P}
\left(
\bigcup_{m=1}^{\infty}C_m
\right)\\
&=
\lim_{m\to\infty}\mathbb{P}(C_m)\\
&=0.
\end{aligned}
\]

Since
\[
\{A_n\text{ i.o.}\}^c
=
\{A_n\text{ only finitely often}\},
\]
we obtain
\[
\mathbb{P}(A_n\text{ i.o.})
=
1-
\mathbb{P}(A_n\text{ only finitely often})
=
1.
\]

Thus
\[
\boxed{
\mathbb{P}(A_n\text{ i.o.})=1.
}
\]

\end{proof}


%============================================================
% Comparison of the Two Lemmas
%============================================================

\subsection{Comparison of the Two Borel--Cantelli Lemmas}

The two Borel--Cantelli lemmas can be summarized as follows.

\begin{theorem}{Borel--Cantelli Summary}
{borel-cantelli-summary}

For a sequence of events $(A_n)_{n\geq1}$,

\[
\sum_{n=1}^{\infty}\mathbb{P}(A_n)<\infty
\quad\Longrightarrow\quad
\mathbb{P}(A_n\text{ i.o.})=0.
\]

If, in addition, the events are independent, then
\[
\sum_{n=1}^{\infty}\mathbb{P}(A_n)=\infty
\quad\Longrightarrow\quad
\mathbb{P}(A_n\text{ i.o.})=1.
\]

\end{theorem}

\begin{observation}

Independence is \emph{not} required for the first Borel--Cantelli
lemma, whereas independence is required for the second Borel--Cantelli
lemma.

\end{observation}


%============================================================
% Strong Law with Fourth Moments
%============================================================

\section{Strong Law Under a Fourth-Moment Assumption}

\subsection{Statement of the Result}

\begin{theorem}{Strong Law with Finite Fourth Moment}
{strong-law-fourth-moment}

Let
\(
X_1,X_2,\ldots
\)
be independent and identically distributed random variables satisfying
\(
\mathbb{E}[X_1]=\mu
\)
and
\[
\mathbb{E}[X_1^4]<\infty.
\]

If
\[
S_n=X_1+\cdots+X_n,
\]
then
\[
\boxed{
\frac{S_n}{n}
\xrightarrow{\mathrm{a.s.}}
\mu.
}
\]

\end{theorem}


\subsection{Proof}

\begin{proof}

By replacing $X_i$ by $X_i-\mu$, it is enough to prove the result
when
\[
\mu=0.
\]

Thus assume
\[
\mathbb{E}[X_i]=0
\]
and
\[
\mathbb{E}[X_i^4]<\infty.
\]

We shall prove that
\[
\frac{S_n}{n}
\xrightarrow{\mathrm{a.s.}}
0.
\]

Consider the fourth moment of $S_n$:
\[
\mathbb{E}(S_n^4)
=
\mathbb{E}
\left[
\left(
\sum_{i=1}^{n}X_i
\right)^4
\right].
\]

Expanding the fourth power gives
\[
S_n^4
=
\sum_{i_1,i_2,i_3,i_4=1}^{n}
X_{i_1}X_{i_2}X_{i_3}X_{i_4}.
\]

Because the $X_i$ are independent and
\[
\mathbb{E}[X_i]=0,
\]
any term containing an index that appears only once has expectation
zero.

The only nonzero contributions are therefore of the following two
types:

\begin{enumerate}
    \item all four indices are equal;
    \item two distinct indices each occur twice.
\end{enumerate}

For the first type,
\[
\sum_{i=1}^{n}
\mathbb{E}[X_i^4]
=
n\mathbb{E}[X_1^4].
\]

For the second type, choose two distinct indices $i\neq j$.
The corresponding term is
\[
X_i^2X_j^2.
\]

For each pair $\{i,j\}$ there are
\[
\frac{4!}{2!2!}=6
\]
ways to arrange the indices. Hence the contribution is
\[
6
\sum_{1\leq i<j\leq n}
\mathbb{E}[X_i^2X_j^2].
\]

By independence,
\[
\mathbb{E}[X_i^2X_j^2]
=
\mathbb{E}[X_i^2]
\mathbb{E}[X_j^2].
\]

Since the variables are identically distributed,
\[
\mathbb{E}[X_i^2X_j^2]
=
\left(
\mathbb{E}[X_1^2]
\right)^2.
\]

Therefore,
\[
\begin{aligned}
\mathbb{E}[S_n^4]
&=
n\mathbb{E}[X_1^4]
+
6\binom{n}{2}
\left(
\mathbb{E}[X_1^2]
\right)^2\\
&=
n\mathbb{E}[X_1^4]
+
3n(n-1)
\left(
\mathbb{E}[X_1^2]
\right)^2.
\end{aligned}
\]

Hence there exists a constant $C<\infty$ such that
\[
\boxed{
\mathbb{E}[S_n^4]\leq Cn^2.
}
\]

\medskip

Now apply Markov's inequality to the nonnegative random variable
$S_n^4$. For every $\varepsilon>0$,
\[
\begin{aligned}
\mathbb{P}
\left(
\left|
\frac{S_n}{n}
\right|>\varepsilon
\right)
&=
\mathbb{P}
\left(
|S_n|>\varepsilon n
\right)\\
&=
\mathbb{P}
\left(
S_n^4>\varepsilon^4n^4
\right)\\
&\leq
\frac{\mathbb{E}[S_n^4]}
{\varepsilon^4n^4}.
\end{aligned}
\]

Using
\[
\mathbb{E}[S_n^4]\leq Cn^2,
\]
we obtain
\[
\mathbb{P}
\left(
\left|
\frac{S_n}{n}
\right|>\varepsilon
\right)
\leq
\frac{C}{\varepsilon^4n^2}.
\]

Therefore,
\[
\sum_{n=1}^{\infty}
\mathbb{P}
\left(
\left|
\frac{S_n}{n}
\right|>\varepsilon
\right)
\leq
\frac{C}{\varepsilon^4}
\sum_{n=1}^{\infty}\frac1{n^2}
<\infty.
\]

By the first Borel--Cantelli lemma,
\[
\mathbb{P}
\left(
\left|
\frac{S_n}{n}
\right|>\varepsilon
\text{ i.o.}
\right)
=0.
\]

Thus, for every $\varepsilon>0$,
\[
\left|
\frac{S_n}{n}
\right|
\leq\varepsilon
\qquad
\text{eventually almost surely}.
\]

Hence
\[
\frac{S_n}{n}
\xrightarrow{\mathrm{a.s.}}
0.
\]

Returning to the original variables, write
\[
X_i=(X_i-\mu)+\mu.
\]

Then
\[
\frac{S_n}{n}
=
\frac1n
\sum_{i=1}^{n}(X_i-\mu)
+\mu.
\]

The first term converges almost surely to $0$. Therefore,
\[
\boxed{
\frac{S_n}{n}
\xrightarrow{\mathrm{a.s.}}
\mu.
}
\]

\end{proof}





\lectureend


%%%%%%%%%%%%%%%%%%%%%%%%%%%%%%%%%%%%%%%%%%%%%%%%%%%%%%%%%%%%
%% Quiz 7
%%%%%%%%%%%%%%%%%%%%%%%%%%%%%%%%%%%%%%%%%%%%%%%%%%%%%%%%%%%%

\begin{quiz}

Let
\(
X_1,X_2,\ldots
\)
be i.i.d.\ real-valued random variables and let
\(
S_n=X_1+\cdots+X_n.
\)

Suppose that there is a sequence of constants
\(
a_n\in\mathbb{R}
\)
such that
\[
\frac{S_n}{n}-a_n
\xrightarrow{P}
0.
\]

Prove that
\[
\boxed{
x\mathbb{P}(|X_1|>x)
\longrightarrow0
\qquad\text{as }x\to\infty.
}
\]

\begin{enumerate}

    %--------------------------------------------------------
    % Part (a)
    %--------------------------------------------------------

    \item[(a)]
    Let
    \(
    X_1',X_2',\ldots
    \)
    be an independent copy of the sequence $(X_i)_{i\geq1}$ and set
    \[
    Y_i=X_i-X_i'.
    \]

    Show that the $Y_i$'s are i.i.d.\ and symmetric, and that
    \[
    \frac1n\sum_{i=1}^nY_i
    \xrightarrow{P}
    0.
    \]

    %--------------------------------------------------------
    % Part (b)
    %--------------------------------------------------------

    \item[(b)] \textbf{A maximal inequality.}

    Let
    \[
    T_k=Y_1+\cdots+Y_k,
    \qquad
    1\leq k\leq n.
    \]

    Fix $t>0$.

    \begin{enumerate}

        \item[(i)]
        Define the first crossing time
        \[
        \tau
        =
        \min\{k:T_k>t\},
        \]
        with $\tau=\infty$ if the set is empty.

        On the event $\{\tau\leq n\}$, consider the transformation
        which reflects all the increments after the first crossing:
        \[
        (Y_1,\ldots,Y_\tau,Y_{\tau+1},\ldots,Y_n)
        \longmapsto
        (Y_1,\ldots,Y_\tau,-Y_{\tau+1},\ldots,-Y_n).
        \]

        Explain why this transformation preserves probabilities.

        \item[(ii)]
        If $\widetilde{T}_n$ denotes the endpoint after reflection,
        show that
        \[
        \widetilde{T}_n
        =
        2T_\tau-T_n.
        \]

        Deduce that
        \[
        \mathbb{P}
        \left(
        \max_{1\leq k\leq n}T_k>t
        \right)
        \leq
        2\mathbb{P}(T_n>t).
        \]

        \item[(iii)]
        Apply the same argument to $-T_k$ to show that
        \[
        \mathbb{P}
        \left(
        \max_{1\leq k\leq n}|T_k|>t
        \right)
        \leq
        2\mathbb{P}(|T_n|>t).
        \]

        \item[(iv)]
        Show that
        \[
        \max_{1\leq i\leq n}|Y_i|>2t
        \quad\Longrightarrow\quad
        \max_{1\leq k\leq n}|T_k|>t.
        \]

        Conclude that
        \[
        \mathbb{P}
        \left(
        \max_{1\leq i\leq n}|Y_i|>2t
        \right)
        \leq
        2\mathbb{P}
        \left(
        \left|
        \sum_{i=1}^nY_i
        \right|
        >t
        \right).
        \]

    \end{enumerate}

    %--------------------------------------------------------
    % Part (c)
    %--------------------------------------------------------

    \item[(c)]
    Put
    \[
    p_n
    =
    \mathbb{P}(|Y_1|>\varepsilon n).
    \]

    Show that
    \[
    \mathbb{P}
    \left(
    \max_{1\leq i\leq n}|Y_i|>\varepsilon n
    \right)
    =
    1-(1-p_n)^n.
    \]

    Deduce that
    \[
    n\mathbb{P}
    \left(
    |X_1-X_1'|>\varepsilon n
    \right)
    \longrightarrow0.
    \]

    %--------------------------------------------------------
    % Part (d)
    %--------------------------------------------------------

    \item[(d)]
    Show that there exists $M<\infty$ such that
    \[
    \mathbb{P}(|X_1'|\leq M)\geq\frac12.
    \]

    Show that, for $x>2M$,
    \[
    \{
    |X_1|>x,\ |X_1'|\leq M
    \}
    \subseteq
    \{
    |X_1-X_1'|>x/2
    \}.
    \]

    Hence prove that
    \[
    \leq
    2\mathbb{P}
    \left(
    |X_1-X_1'|>\frac{x}{2}
    \right),
    \qquad x>2M.
    \]

    %--------------------------------------------------------
    % Part (e)
    %--------------------------------------------------------

    \item[(e)]
    Combine the preceding parts to conclude that
    \[
    x\mathbb{P}(|X_1|>x)
    \longrightarrow0
    \qquad\text{as }x\to\infty.
    \]

\end{enumerate}

\end{quiz}





